Managed agent platforms let one cloud account run many agents, each with an
IAM role of its own. Tools that audit IAM fix one identity and ask what it can
reach. That view cannot express the question this paper is about: whether one
agent can take over another agent in the same account and act with the second
agent's permissions.

In a live AWS account, an attacker principal holding one AWS managed policy
and nothing else replaces a sibling agent's definition, runs that code under
the victim's execution role, and reads back a value it cannot read directly.
The gate is a name match in that policy's \iam{iam:PassRole} statement, which
the console's own default role names satisfy, so what stands between two
agents is a check on the kind of role rather than on whose role it is: a
control role built identically but named differently is refused. The takeover
also outlives the attacker. With every policy detached from the attacker
principal, the operator's own call still returns the planted code's answer. A
published route reaches the same role more cheaply and leaves the agent as it
was; this one leaves the agent answering its own callers with attacker code.

We state the escalation as a predicate over what a principal reaches in effect
and what it was authorised to reach, and implement it in \tool, which reads
infrastructure as code and reports a cross-agent escalation together with the
permission that is missing from it. Over \nProjects public Terraform projects
that deploy agent runtimes and the CloudFormation and CDK deployments in the
same collection, \nSharedRoleRuntimes of \nMultiAgentRuntimes multi-agent
deployments run every agent under one execution role, which makes the
escalation unnecessary rather than impossible, and \nManagedPolicyConfigs
attach the managed policy the live experiment uses. A further \nBlocked
takeover candidates are blocked on \iam{iam:PassRole} alone,
\nBlockedAgentToAgent of them between two agents of the same deployment. We
find \nChainsFound escalations that are complete in source, an upper bound of
\chainBound. A \foundry corpus of \nFoundryDirs template directories is
surveyed separately, because agents there are not declared at deployment time.
